\section{はじめに}
クリフディメンションは，上下に動く突起（クリフ A）に両手でぶら下がった状態から体を振って離手し，
空中で向きを変え，左右に動く突起（クリフ B）を両手で捕捉する課題である（図~\ref{fig:device3d}）．
突起の奥行きは 3\,cm しかなく，成否は重心の放物線だけでなく，手先の相対速度，掛かり代，指の把持容量に左右される．
本研究の最終目標は，身長・体重・筋力の異なる誰にとっても「どう飛べば届くか」を計算し，
それを学習器が自律的に獲得できるようにすることである．

この課題は学習にとって難しい．捕捉は 3\,cm の縁に手が掛かったかどうかで決まるので報酬が極端に疎であり，
成功軌道が一度も観測されないまま学習が停滞する（第~\ref{sec:rl}節で実測する）．
そこで本稿は三段階で進める．
第一段（第~\ref{sec:model}--\ref{sec:results}節）は，学習に先立ち「物理的に何が要求されるか」を決定論的な軌道最適化で確定する．
第二段（第~\ref{sec:dfl}節）は，その最適化が副産物として与える\textbf{双対変数}（価値の勾配と制約の影の価格）を
学習の教師にする双対場学習を提案する．
第三段（第~\ref{sec:protocol}--\ref{sec:learning}節）は，計画した運動が実行できるか，学習した場が制御に使えるかを，
評価規約を固定した 10 個の実験で検証する．

本稿の貢献と主な結果は次のとおりである．
(i)~前腕と突起の干渉を含む平面 5 リンク人体モデルの多相軌道最適化により，装置周期 9 通り，体重 6 通り，離手位相 16 通りの
\gridN 条件について必要把持容量を求めた．必要容量は体重で決まり装置周期にほとんど依らないこと，
捕捉後に体を支える保持が必要容量を決めていることを，乗数と再最適化による介入の両方で示す．
(ii)~最適化の双対変数を学習対象とする双対場学習を定式化し，双対ラベルの検査，場の適合，候補の順位付け，閉ループ制御，
データの量と集め方，体格への汎化を，双対ラベルを使わない同一構成の制御器と対にして比較した．
(iii)~最適化器と同じ記号力学を共有する高速な環境を作り，計画の再生が失敗する原因を制御周期，離手時刻，計画メッシュ，
開始相に分解した．
結論を先に述べると，双対ラベルは価値の勾配の精度，0.8\,s 先の候補の順位付け，離手に至る割合を改善するが，
成功率は双対ラベルの有無にかかわらず \clSmpcSuccInt\,\%であった（差 $\Delta S=\deltaS$ ポイント）．
参照解の再生でも成功率は \clReplaySuccInt\,\%にとどまり，この課題の難しさは振りの学習よりも，
活性制約の上を通る最適解を実行することにある．本稿はこの否定的な結果を，原因の分解とともに報告する．
